% =========================================================
% Chapter 1 — Basic statistical concepts
% Time Series Analysis — Aymen Ben Brik (aymen.benbrik@esprit.tn)
% Esprit School of Business
%
% Chapter content (no preamble), to be \input by polycopie.tex or
% chapitre1-statistics/chapitre-standalone.tex.
% =========================================================

\chapter{Basic statistical concepts}
\label{ch:basic-stats}

\section*{Introduction}
\addcontentsline{toc}{section}{Introduction}

Before studying time series proper, this chapter reviews the
statistical concepts on which the rest of the course relies:
\emph{moments} of a distribution, classical \emph{point estimators},
\emph{Method of Moments} and \emph{Maximum Likelihood Estimation},
sampling distributions, and basic \emph{inference} (confidence
intervals and hypothesis tests).

\medskip
The convention throughout the chapter is $X_1, \dots, X_n$ are
\emph{independent and identically distributed} (i.i.d.) draws from a
distribution with density $f(x;\theta)$ parameterised by some unknown
$\theta$. We observe a sample and want to learn $\theta$. Time series
will later violate the i.i.d.\ assumption, but the basic estimators
introduced here remain the building blocks.

% =========================================================
\section{Moments of a distribution}
\label{sec:moments}

\begin{defbox}[Raw and central moments]
\label{def:moments}
Let $X$ be a real-valued random variable with density $f$ and mean
$\mu = \E[X]$.
\begin{itemize}
  \item The \emph{raw $\ell$-th moment} is
        $m'_\ell \;=\; \E[X^\ell] \;=\; \int x^\ell f(x)\, dx$.
  \item The \emph{central $\ell$-th moment} is
        $m_\ell \;=\; \E[(X - \mu)^\ell]$.
\end{itemize}
\end{defbox}

\paragraph{Useful quantities.}
\begin{align*}
  \mu          &= m'_1, &
  \sigma^2     &= m_2 = \Var[X], \\[2pt]
  \text{skewness } S(X) &= \frac{m_3}{\sigma^3}, &
  \text{kurtosis } K(X) &= \frac{m_4}{\sigma^4}.
\end{align*}

\begin{rembox}[Interpretation]
\begin{itemize}
  \item $S(X)$ measures \emph{asymmetry}. $S = 0$ for any symmetric
        distribution; $S > 0$ for a right-skewed distribution.
  \item $K(X)$ measures \emph{tail-heaviness}. The Gaussian has $K = 3$;
        the \emph{excess kurtosis} $K - 3$ is therefore the standard
        comparison: $K - 3 > 0$ indicates heavier tails than the
        Gaussian (leptokurtic).
\end{itemize}
\end{rembox}

\begin{exbox}[Moments of the standard normal]
For $X \sim \mathcal{N}(\mu, \sigma^2)$ one shows by direct integration
or via the moment-generating function:
\[
  \E[X] = \mu, \quad \Var[X] = \sigma^2, \quad m_3 = 0, \quad m_4 = 3\,\sigma^4.
\]
Hence $S(X) = 0$ and $K(X) = 3$.
\end{exbox}

% =========================================================
\section{Classical estimators and sampling distributions}
\label{sec:estimators}

\subsection{Sample mean and sample variance}

\begin{defbox}[Empirical estimators]
Given i.i.d.\ data $X_1, \dots, X_n$:
\[
  \bar{X}_n \;=\; \frac{1}{n}\sum_{i=1}^{n} X_i,
  \qquad
  S_n^2 \;=\; \frac{1}{n-1}\sum_{i=1}^{n}(X_i - \bar{X}_n)^2.
\]
\end{defbox}

\begin{propbox}[Unbiasedness of $\bar{X}_n$ and $S_n^2$]
\label{prop:unbiased}
For i.i.d.\ data with finite second moments, $\E[\bar{X}_n] = \mu$ and
$\E[S_n^2] = \sigma^2$.
\end{propbox}

\begin{proof}
For the mean: $\E[\bar{X}_n] = \tfrac{1}{n}\sum_i \E[X_i] = \mu$ by
linearity. For the variance, expand:
\[
  \sum_i (X_i - \bar{X}_n)^2 = \sum_i X_i^2 - n\, \bar{X}_n^2.
\]
Take expectations: $\E[X_i^2] = \mu^2 + \sigma^2$ and
$\E[\bar{X}_n^2] = \mu^2 + \sigma^2 / n$. Hence
$\E[\sum_i(X_i - \bar{X}_n)^2] = n(\mu^2 + \sigma^2) - n(\mu^2 + \sigma^2/n)
= (n-1)\,\sigma^2$. The $n - 1$ denominator in $S_n^2$ ensures
$\E[S_n^2] = \sigma^2$.
\end{proof}

\begin{rembox}[Why $n - 1$?]
Replacing the unknown $\mu$ by $\bar{X}_n$ in the variance formula
``costs'' one degree of freedom. Dividing by $n$ would systematically
underestimate $\sigma^2$.
\end{rembox}

\subsection{Sampling distribution under normality}

\begin{thbox}[Cochran-type results]
\label{th:sampling-normal}
If $X_1, \dots, X_n \iid \mathcal{N}(\mu, \sigma^2)$, then:
\begin{enumerate}
  \item $\bar{X}_n \sim \mathcal{N}(\mu,\, \sigma^2 / n)$;
  \item $\dfrac{(n-1)\, S_n^2}{\sigma^2} \sim \chi^2_{n-1}$;
  \item $\bar{X}_n$ and $S_n^2$ are independent;
  \item $\dfrac{\bar{X}_n - \mu}{S_n / \sqrt{n}} \sim t_{n-1}$.
\end{enumerate}
\end{thbox}

\begin{rembox}
Result (4) follows from (1)--(3): the numerator is standard normal, the
denominator is the square root of an independent $\chi^2_{n-1}$ divided
by its degrees of freedom — the very definition of Student's $t$.
\end{rembox}

\subsection{Asymptotic results}

\begin{thbox}[Law of Large Numbers]
\label{th:lln}
For i.i.d.\ data with $\E[X_i] = \mu$:
\[
  \bar{X}_n \;\xrightarrow{p}\; \mu \quad \text{as } n \to \infty.
\]
\end{thbox}

\begin{thbox}[Central Limit Theorem]
\label{th:clt}
For i.i.d.\ data with $\E[X_i] = \mu$ and $\Var[X_i] = \sigma^2 < \infty$:
\[
  \frac{\bar{X}_n - \mu}{\sigma / \sqrt{n}}
  \;\xrightarrow{d}\;
  \mathcal{N}(0, 1) \quad \text{as } n \to \infty.
\]
\end{thbox}

\begin{rembox}[Practical consequence]
Even when the $X_i$ are not Gaussian, the sample mean behaves
approximately as a normal random variable for large $n$. The rule of
thumb $n \geq 30$ is conservative; in practice the convergence is
faster for symmetric, light-tailed distributions and slower for skewed
or heavy-tailed ones.
\end{rembox}

% =========================================================
\section{Estimation methods}
\label{sec:estimation}

\subsection{Method of Moments (MoM)}

\begin{defbox}[Method of Moments]
\label{def:mom}
Let $\theta \in \R^k$. The MoM estimator $\widehat{\theta}_n^{\text{MoM}}$
solves the system of equations
\[
  \E_\theta[X^p] \;=\; \tfrac{1}{n}\sum_{i=1}^{n} X_i^p,
  \qquad p = 1, 2, \dots, k.
\]
That is, theoretical moments are equated to sample moments.
\end{defbox}

\begin{exbox}[$\mathcal{N}(\mu, \sigma^2)$ via MoM]
With $\theta = (\mu, \sigma^2)$ and $k = 2$:
\begin{align*}
  \mu        &= \bar{X}_n, \\
  \mu^2 + \sigma^2 &= \tfrac{1}{n}\sum_i X_i^2.
\end{align*}
Subtracting yields
$\widehat{\sigma}^2_{\text{MoM}} = \tfrac{1}{n}\sum(X_i - \bar{X}_n)^2$,
which differs from $S_n^2$ by a factor $n / (n - 1)$.
\end{exbox}

\subsection{Maximum Likelihood Estimation (MLE)}

\begin{defbox}[Likelihood and MLE]
For i.i.d.\ data with density $f(x; \theta)$:
\[
  L(\theta) \;=\; \prod_{i=1}^{n} f(X_i; \theta),
  \qquad
  \ell(\theta) = \log L(\theta).
\]
The MLE is $\widehat{\theta}_n^{\text{MLE}}
\;=\; \arg\max_\theta\, \ell(\theta)$.
\end{defbox}

\begin{exbox}[$\mathcal{N}(\mu, \sigma^2)$ via MLE]
\label{ex:mle-gaussian}
The log-likelihood is
\[
  \ell(\mu, \sigma^2)
  = -\tfrac{n}{2}\log(2\pi\sigma^2)
    - \tfrac{1}{2\sigma^2}\sum_{i=1}^{n}(X_i - \mu)^2.
\]
Setting $\partial \ell / \partial \mu = 0$ gives
$\widehat{\mu}_{\text{MLE}} = \bar{X}_n$. Then setting
$\partial \ell / \partial \sigma^2 = 0$ gives
$\widehat{\sigma}^2_{\text{MLE}}
= \tfrac{1}{n}\sum(X_i - \bar{X}_n)^2$. So MLE coincides with MoM here,
and both are biased for $\sigma^2$ (their expectation is
$\tfrac{n-1}{n}\sigma^2$).
\end{exbox}

\begin{thbox}[Asymptotic normality of the MLE]
\label{th:mle-asymp}
Under regularity conditions (identifiability, smoothness, finite Fisher
information $I(\theta_0) > 0$):
\[
  \sqrt{n}\,\bigl(\widehat{\theta}_n^{\text{MLE}} - \theta_0\bigr)
  \;\xrightarrow{d}\;
  \mathcal{N}\!\bigl(0,\, I(\theta_0)^{-1}\bigr).
\]
\end{thbox}

\begin{rembox}[MoM vs.\ MLE]
\begin{itemize}
  \item MoM is computationally simple — just match moments.
  \item MLE is \emph{asymptotically efficient}: its variance attains
        the Cramér--Rao lower bound $I(\theta_0)^{-1}/n$.
  \item Both estimators are consistent (under regularity).
  \item For ``nice'' families (Gaussian, exponential, Poisson), MoM
        and MLE often coincide or are very close.
\end{itemize}
\end{rembox}

% =========================================================
\section{Multivariate data}
\label{sec:multivariate}

\subsection{Joint, marginal, and conditional distributions}

\begin{defbox}[Joint and marginal density]
For two random variables $X, Y$ with joint density $f(x, y)$:
\[
  f_X(x) \;=\; \int f(x, y)\, dy,
  \qquad
  f_{Y \mid X}(y \mid x) \;=\; \frac{f(x, y)}{f_X(x)}.
\]
The chain rule reads $f(x, y) = f_{Y \mid X}(y \mid x)\, f_X(x)$.
\end{defbox}

\begin{defbox}[Independence]
$X$ and $Y$ are \emph{independent} iff $f(x, y) = f_X(x)\, f_Y(y)$ for
all $x, y$.
\end{defbox}

\subsection{Covariance and correlation}

\begin{defbox}[Covariance and correlation]
\[
  \Cov(X, Y) = \E\bigl[(X - \E X)(Y - \E Y)\bigr],
  \qquad
  \Corr(X, Y) = \frac{\Cov(X, Y)}{\sqrt{\Var X \Var Y}}.
\]
\end{defbox}

\begin{propbox}[Sample versions]
\[
  \widehat{\Cov}(X, Y) = \tfrac{1}{n - 1}\sum_{i=1}^{n}
    (X_i - \bar X)(Y_i - \bar Y),
  \qquad
  \widehat{\Corr}(X, Y) = \frac{\widehat{\Cov}(X, Y)}{\widehat{\sigma}_X\,\widehat{\sigma}_Y}.
\]
The sample correlation lies in $[-1, 1]$.
\end{propbox}

\begin{warnbox}[Correlation captures only linear dependence]
$\Corr(X, Y) = 0$ does not imply independence. Counter-example: let
$X$ be symmetric around $0$ and $Y = X^2$. Then $\E[XY] = \E[X^3] = 0
= \E[X]\E[Y]$, hence $\Cov(X, Y) = 0$, yet $Y$ is a deterministic
function of $X$.
\end{warnbox}

% =========================================================
\section{Statistical inference}
\label{sec:inference}

\subsection{Confidence intervals for the mean}

\begin{defbox}[Confidence interval]
A $(1 - \alpha)$-confidence interval for $\mu$ is a random interval
$[\widehat{L}_n, \widehat{U}_n]$ such that
$\Prob\bigl(\widehat{L}_n \leq \mu \leq \widehat{U}_n\bigr) = 1 - \alpha$
in repeated sampling.
\end{defbox}

\paragraph{Gaussian sample, known $\sigma$ ($z$-interval).}
Let $X_1, \dots, X_n \iid \mathcal{N}(\mu, \sigma^2)$ with $\sigma$
known. By Theorem~\ref{th:sampling-normal} (1):
\[
  \mathrm{CI}_{1-\alpha}(\mu)
  = \left[\,\bar{X}_n - z_{1-\alpha/2}\,\tfrac{\sigma}{\sqrt{n}},\;
            \bar{X}_n + z_{1-\alpha/2}\,\tfrac{\sigma}{\sqrt{n}}\,\right].
\]

\paragraph{Gaussian sample, unknown $\sigma$ ($t$-interval).}
Replace $\sigma$ by $S_n$ and use Student's $t$ (Theorem~\ref{th:sampling-normal} (4)):
\[
  \mathrm{CI}_{1-\alpha}(\mu)
  = \left[\,\bar{X}_n - t_{n-1, 1-\alpha/2}\,\tfrac{S_n}{\sqrt{n}},\;
            \bar{X}_n + t_{n-1, 1-\alpha/2}\,\tfrac{S_n}{\sqrt{n}}\,\right].
\]

\subsection{Hypothesis testing}

\begin{defbox}[Test on the mean]
Two-sided test:
$H_0: \mu = \mu_0$ vs.\ $H_1: \mu \neq \mu_0$. Test statistic
\[
  T \;=\; \frac{\bar{X}_n - \mu_0}{S_n / \sqrt{n}}
  \;\sim\; t_{n-1} \text{ under } H_0.
\]
Reject $H_0$ at level $\alpha$ if $|T| > t_{n-1, 1-\alpha/2}$, i.e.\ the
$p$-value $2 \cdot \Prob(T_{n-1} > |T|) < \alpha$.
\end{defbox}

\begin{rembox}[Errors]
\begin{itemize}
  \item Type I error (size of the test): $\alpha = \Prob(\text{reject } H_0 \mid H_0)$.
  \item Type II error: $\beta = \Prob(\text{not reject } H_0 \mid H_1)$.
  \item \emph{Power} of the test: $1 - \beta$, the probability of
        correctly rejecting $H_0$ when $H_1$ is true.
\end{itemize}
\end{rembox}

\begin{exbox}[Numerical illustration]
Suppose $n = 50$ Atlanta monthly average temperatures with sample
mean $\bar{X}_{50} = 61.2$ and sample standard deviation $S_{50} = 16.8$.
Test $H_0 : \mu = 60$ against $H_1 : \mu \neq 60$ at $\alpha = 0.05$.
\[
  T \;=\; \frac{61.2 - 60}{16.8 / \sqrt{50}}
       \;\approx\; 0.505,
  \qquad
  t_{49, 0.975} \;\approx\; 2.010.
\]
Since $|T| < 2.010$, we do \emph{not} reject $H_0$;
$p\text{-value} \approx 2 \cdot \Prob(t_{49} > 0.505) \approx 0.616$.
\medskip
The corresponding 95\% CI is
$[61.2 - 4.78,\, 61.2 + 4.78] = [56.4,\, 66.0]$, which contains
$\mu_0 = 60$, consistent with the test conclusion.
\end{exbox}

\begin{warnbox}[Common pitfalls]
\begin{itemize}
  \item ``Failing to reject $H_0$'' is \emph{not} ``proving $H_0$''.
  \item A small $p$-value does not imply a large effect — it only
        signals statistical significance, not practical importance.
  \item Multiple testing inflates Type I error: use Bonferroni or FDR
        corrections when running many tests.
\end{itemize}
\end{warnbox}

% =========================================================
\section*{Summary}
\addcontentsline{toc}{section}{Summary}

\begin{center}
\begin{tabular}{lll}
\toprule
\textbf{Quantity} & \textbf{Estimator} & \textbf{Sampling distribution (Gaussian data)} \\
\midrule
$\mu$         & $\bar{X}_n = \tfrac{1}{n}\sum X_i$
              & $\mathcal{N}(\mu, \sigma^2/n)$ \\
$\sigma^2$    & $S_n^2 = \tfrac{1}{n-1}\sum (X_i - \bar X_n)^2$
              & $(n-1)S_n^2/\sigma^2 \sim \chi^2_{n-1}$ \\
Pivot         & $T = (\bar{X}_n - \mu)/(S_n/\sqrt{n})$
              & $t_{n-1}$ \\
\bottomrule
\end{tabular}
\end{center}

\medskip
\noindent
\textbf{What to remember.}
\begin{itemize}
  \item Two estimation routes: MoM (match moments) and MLE
        (maximise likelihood). Both consistent; MLE asymptotically
        efficient.
  \item Under normality, the sampling distribution of $\bar X$ is
        exactly Gaussian; the CLT extends this asymptotically.
  \item Confidence intervals and hypothesis tests on the mean rely on
        Student's $t$ when $\sigma$ is unknown.
\end{itemize}
